%======================================================================
\section{Introduction}
%======================================================================

\subsection{The linear problem behind explicit-formula conductor bounds}

Weil's explicit formula \cite{Wei52} (in the form of Guinand \cite{Gui48} for $\zeta$) is a linear identity between two
measures: one on the zeros of $\zeta$, one on the prime powers. For even test functions $F$ that are analytic and decay in a
strip, and with $\hF(\xi)=\int F(t)e^{-2\pi i\xi t}\dd t$, it reads
\begin{equation}\label{eq:intro-EF}
F\big(\tfrac i2\big)+F\big(-\tfrac i2\big)+\frac1{2\pi}\int_\RR F(t)\Big(\Re\psi\big(\tfrac14+\tfrac{it}2\big)-\log\pi\Big)\dd t
=\sum_\rho F\Big(\frac{\rho-\frac12}i\Big)+\frac1\pi\sum_{n\ge2}\frac{\Lambda(n)}{\sqrt n}\,\hF\Big(\frac{\log n}{2\pi}\Big).
\end{equation}
The left side, which we denote by $\Arch(F)$, contains only the archimedean data of $\zeta$: the Gamma factor and the pole at
$s=1$. Under RH every zero contributes $F(\gamma)$ with $\gamma$ real, so the right side is $\int F\dd\mu+\frac1\pi\int\hF\dd\nu$ for two
positive measures: $\mu$ on the real ordinates of the zeros and $\nu$ on the points $\log n/(2\pi)$, $n\ge2$.

Positivity of this kind drives the explicit-formula method for lower bounds of discriminants and conductors: the analytic method of
Stark and Odlyzko, recast by Serre through the explicit formula and refined by Poitou \cite{Sta74,Odl76,Odl77,Ser75,Poi77,Poi77b} (see
\cite[pp.~120--121]{Odl90}). A test function with $F\ge0$ on $\RR$ and $\hF\ge0$ on the relevant range makes
the zero side and the prime side non-negative, and \eqref{eq:intro-EF} then bounds the archimedean data, in particular the
conductor. Odlyzko asked which test functions give the best bounds \cite[Open Problem~2.2, p.~123]{Odl90}, how the prime ideals and the
zeros share the explicit formula for minimal discriminants \cite[Open Problem~6.1 and Table~5]{Odl90}, and whether the bounds
obtained under GRH remain valid without it \cite[Open Problem~6.3, p.~132]{Odl90}. For $n=1$ his tables of root-discriminant
bounds, that is, for data of the type of $\zeta$, give $0.997$, both under GRH and unconditionally \cite[Table~4, p.~134]{Odl90};
these values are based on his 1976 tables \cite{Odl76} (which print $0.996$), optimised over a one-parameter family of kernels. The true conductor is $1$.

This paper studies the linear problem behind these bounds for the data of $\zeta$ itself, as a problem about \emph{all}
positive solutions of \eqref{eq:intro-EF}. Its main idea is that the two questions it raises, how sharp the optimal
conductor bound is and how many positive solutions there are, are linked by complementary slackness: an extremal test function for
the conductor bound pins down the support of every positive solution (Lemma~\ref{lem:compslack}). Whether, conversely, uniqueness
forces such an extremal function is open (Problem~\ref{op:strict-compl}).

\subsection{Admissible pairs, the slack, and Conjectures U and S}\label{ssec:intro-pairs}

An \emph{admissible pair} (Definition~\ref{def:admissible}) is a pair $(\mu,\nu)$ of positive measures, $\mu$ even on $\RR$ and
$\nu$ on $[\xin2,\infty)$, where $\xin2=\log2/(2\pi)$, such that $\Arch(F)=\int F\dd\mu+\frac1\pi\int\hF\dd\nu$ for every $F$ in a natural
test class $\Tc$ (Definition~\ref{def:testclass}). Nothing is assumed about where $\mu$ and $\nu$ live, beyond the gap
$(-\xin2,\xin2)$ in the prime measure, and nothing about integrality. We write $\Kad$ for the set of admissible pairs and $\pz=(\muz,\nuz)$
for the pair of $\zeta$: $\muz$ counts the real zero ordinates with multiplicity and $\nuz=\sum_n\Lambda(n)n^{-1/2}\delta_{\log n/2\pi}$. The
pair $\pz$ is admissible if and only if RH holds. We consider three statements:
\[
\condE:\ \Kad\ne\emptyset,\qquad \condU:\ \Kad\subseteq\{\pz\},\qquad \condS:\ \kstar\le0,
\]
where the \emph{slack} $\kstar$ is the infimum of $\Arch(F)$ over the cone $\Cone$ of $F\in\Tc$ with $F\ge0$ on $\RR$, $\hF\ge0$ on
$[\xin2,\infty)$ and $\int F=1$. In the language of conductors, $\qmin=e^{-2\pi\kstar}$ is the least conductor for which the Gamma
factor $\GR$, one simple pole and a positive prime measure on $[\xin2,\infty)$ admit a positive solution of the explicit formula, and
it is the fully optimised degree-one conductor bound (Proposition~\ref{prop:conductor}).

\begin{lconj}{U}\label{conj:U}
Every admissible pair, if there is one, is $\zeta$'s pair: $\Kad\subseteq\{\pz\}$.
\end{lconj}

\begin{lconj}{S}\label{conj:S}
$\kstar\le0$. Equivalently, $\qmin\ge1$; equivalently, no admissible pair has a zero measure $\mu\ge\lambda\dd t$ with $\lambda>0$.
\end{lconj}

The logic of these statements is (Theorem~\ref{thm:logic})
\[
\mathrm{RH}\Rightarrow\condE,\qquad \condU\Rightarrow\condS,\qquad
\condU\iff\big[\condE\Rightarrow\mathrm{RH}\big]\wedge\big[\mathrm{RH}\Rightarrow\Kad=\{\pz\}\big].
\]
So \condU\ is the conjunction of a prime-free criterion for RH and a uniqueness statement under RH. It does \emph{not} imply
RH: it is consistent with $\neg$RH$\wedge\neg$\condE, and under \condU\ alone we only get RH$\Leftrightarrow$\condE. For the field $\QQ$,
every GRH-type conductor bound given by a test function in $\Cone$ holds without RH exactly when $\Arch\ge0$ on $\Cone$, that is,
under \condE; so \condE\ is the statement that all such bounds of Odlyzko's Open Problem~6.3 hold for $\QQ$, and the first half
of \condU\ is a converse to it. (The tabulated value $0.997<1$ holds for $\QQ$ trivially.) Under RH, \condS\ is equivalent to
$\kstar=0$, that is, to the statement that $\zeta$'s data are exactly critical: they cannot be tightened at all without losing
every positive solution.

\subsection{Main results}

We write $\BRS=\{\log m/(4\pi):m\ge1\}$; it contains the prime-power nodes $\log n/(2\pi)=\log(n^2)/(4\pi)$.

\begin{lthm}{A}[Duality and criticality]\label{thm:intro-A}
\begin{enumerate}[label=\textup{(\alph*)}]
\item $\Kad\ne\emptyset$ if and only if $\Arch(F)\ge0$ for every $F\in\Cone$, if and only if $\kstar\ge0$; and $\qmin=e^{-2\pi\kstar}$.
\item Assume \condE. Then $\kstar=0$ if and only if every admissible $\mu$ is purely atomic, if and only if all admissible $\mu$ are
carried by one discrete set.
\end{enumerate}
\end{lthm}

These are Theorem~\ref{thm:duality}, Proposition~\ref{prop:conductor} and Theorem~\ref{thm:crit}; the logic displayed in
Section~\ref{ssec:intro-pairs} is Theorem~\ref{thm:logic}. In (b) the common discrete set is the real zero set of a \emph{weak magic
function}, a limit of near-optimal test functions. Abstract sign and gap conditions alone do not force rigidity
(Section~\ref{ssec:barrier}); the arithmetic input of Theorem~\ref{thm:intro-B} is the Fourier interpolation basis of Bondarenko, Radchenko and Seip \cite{BRS}, which
is attached to the zeros of $\zeta$ and the points $\BRS$ and does not assume RH.

\begin{lthm}{B}[Uniqueness near $\zeta$'s support]\label{thm:intro-B}
Let $(\mu,\nu)\in\Kad$.
\begin{enumerate}[label=\textup{(\alph*)}]
\item If $\mu$ is carried by $\Zz\cup\{0\}$, then RH holds and $(\mu,\nu)=\pz$. No condition on $\nu$ is needed.
\item Let $E\subset\RR\setminus\Zz$ be countable. If $\mu$ is carried by $\Zz\cup E$, $\nu$ is carried by $\BRS$, and
$\sum_{e\in E,\,e>0}\mu(\{e\})\,|\zeta(\frac12+ie)|\,(1+e)^{-1/2}<\infty$, then RH holds and $(\mu,\nu)=\pz$. This holds whenever $E$ is finite.
\end{enumerate}
\end{lthm}

These are Theorem~\ref{thm:zerosupport} and Theorem~\ref{thm:Gfin} applied to $E\cup(-E)$. Part (a), and the extension of (b) from finite to
countable $E$ with its mean-value argument, are due to Astra (OpenAI), contributed during an independent review of an earlier version of
this paper, and independently verified. Theorem~\ref{thm:Gfin} uses positivity of $\nu$ only at the nodes $\log m/(4\pi)$ with
$m\equiv2\pmod3$, a progression that contains none of $\zeta$'s own prime-power nodes, and positivity of $\mu$ only at $0$; it also allows
the weaker condition $\sum_{e\in E,\,0<e\le T}\mu(\{e\})|\zeta(\frac12+ie)|=o(T^{1/2})$.

\begin{introcor}[Magic-function principle]
Let $F\in\Cone$ with $\Arch(F)=0$. If every real zero of $F$ lies in $\Zz\cup\{0\}$, then \condU\ holds. The same holds if the zeros of $\hF$
in $[\xin2,\infty)$ lie in $\BRS$ and $F$ has only finitely many real zeros outside $\Zz$, or countably many satisfying the summability
condition of Corollary~\ref{cor:magic}(b). If moreover \condE\ holds, then RH holds and $\Kad=\{\pz\}$.
\end{introcor}

This is Corollary~\ref{cor:magic}: complementary slackness (Lemma~\ref{lem:compslack}) puts every admissible pair on the real zeros of $F$ and the
zeros of $\hF$ in $[\xin2,\infty)$, and Theorem~\ref{thm:intro-B} applies. It is the chain that links the two halves of the paper:
near-extremal certificates for the conductor problem, whose limit is an exact magic function with all its real zeros in $\Zz\cup\{0\}$,
would prove \condU\ through Theorem~\ref{thm:intro-B}(a), with no condition on the zeros of its transform. Theorem~\ref{thm:intro-D} below states
this for an explicit family, and Corollary~\ref{cor:certroute} gives a route that needs no limit function.

\begin{introthm}[Integer weights]
Let $q>0$, and let $(\mu,\nu)$ be admissible for the data with Gamma factor $\GR$, one simple pole and conductor $q$. If $\mu$ is purely atomic
with integer weights, then $q\ge1$; if moreover $q=1$, then RH holds and $(\mu,\nu)=\pz$.
\end{introthm}

This is Theorem~\ref{thm:integer}. So \condU\ holds if and only if every admissible $\mu$ is purely atomic with integer weights, while \condS\
holds if and only if every admissible $\mu$ is purely atomic (Theorem~\ref{thm:intro-A}(b) and Corollary~\ref{cor:integer}). The proof uses an
exponential lift of the prime measure, Fej\'er's kernel and Hamburger's theorem.

\begin{lthm}{C}[Near-criticality]\label{thm:intro-C}
\begin{enumerate}[label=\textup{(\alph*)}]
\item $\kstar\le\kOPS\le1.4291572\cdot10^{-1060}$. Hence $\qmin\ge1-8.98\cdot10^{-1060}$, and the same bound holds already in the
classical cone $\ConeOPS$.
\item $\kstar\ge-2.7112\cdot10^{-3}$.
\end{enumerate}
\end{lthm}

These are Theorem~\ref{thm:kappa}, Corollary~\ref{cor:qmin} and Proposition~\ref{prop:lowerbound}; they are
computer-assisted and unconditional. Together they give, unconditionally, $1-8.98\cdot10^{-1060}\le\qmin<1.017182$; only under
\condE, for instance under RH, is $\kstar\ge0$, that is, $\qmin\le1$. So Theorem~\ref{thm:intro-C} sharpens the conductor bound that the
explicit-formula method can give for $\zeta$'s data; the arithmetic conductor of $\zeta$ is $1$ in any case. The certificate in (a) is the
zero-killing function $\Xi^2P_J(t^2)$ with $J=111$, where $\Xi$ is Riemann's $\Xi$-function and $P_J$ is the exact interpolant of
Section~\ref{ssec:hermite}. It vanishes at every zero of $\zeta$, on or off the line, so $\Arch$ of it is a sum over prime powers alone, and
its transform has zeros of order two at the first $J$ prime powers, so that sum starts at the $(J+1)$-st prime power. We certify that these
functions lie in $\Cone$, and even in $\ConeOPS$, for $J\in\{10,60,61,110,111\}$, with no approximation of $P_J$
(Proposition~\ref{prop:exactcone}). Since $P_J\ge1$ on $[0,\infty)$ for these $J$, they also bound every positive solution:
$\int\Xi^2\dd\mu\le3.99\cdot10^{-1060}$ for every admissible pair, so $\mu(I)\le3.99\cdot10^{-1060}/\min_I\Xi^2$ for every compact interval $I$ that
avoids $\Zz$ (Corollary~\ref{cor:nearrigid}). Near-criticality thus forces near-rigidity, at low height. The exponent $1060$ is a
computational budget ($J=111$); the evidence for \condS\ is
Table~\ref{tab:kappa-ladder}, in which the certified bounds keep improving from $J=10$ to $J=111$. A certificate of the classical
(Gaussian--Laguerre) kind shows that every admissible pair, whether or not RH holds, puts mass at most $2.4\cdot10^{-13}$ into each gap
between the first six positive zero ordinates of $\zeta$, shrunk by $0.1$ at both ends, and at most $3.8\cdot10^{-8}$ into listed windows between low prime powers
(Proposition~\ref{prop:lowheight}). With the same gap $\xi_2$, and also at their own natural gaps, the archimedean data of $\QQ(\sqrt5)$ and
$\QQ(\sqrt{-3})$ have positive slack (Proposition~\ref{thm:notcritical}); this indicates that near-criticality is a property of $\zeta$'s data
rather than of the method.

\begin{lthm}{D}[A criterion through the zero-killing family]\label{thm:intro-D}
If the exact Hermite family $F_J=\Xi^2P_J(t^2)$ of Section~\ref{ssec:hermite} satisfies the asymptotic hypotheses of
Theorem~\ref{thm:criterion}, then \condS\ and \condU\ hold, and RH$\Leftrightarrow$\condE.
\end{lthm}

Here $P_J$ is the polynomial with $P_J(0)=1$ and degree at most $2J$ for which the transform of $F_J$ has zeros of order at least two at the first $J$
prime powers; the certificate of Theorem~\ref{thm:intro-C}(a) is the member $J=111$ of this family. The hypotheses ask for uniform
growth bounds, no real limit zeros, and non-negativity of $\liminf_J\FT F_J$ beyond the gap; by Theorem~\ref{thm:intro-B}(a) no strict
positivity of the limiting transform is needed.

\subsection{What is claimed and what is not}\label{ssec:claims}

Theorems~\ref{thm:intro-A} and~\ref{thm:intro-B}, the theorem on integer weights and the Corollary are proved; Theorem~\ref{thm:intro-C} is computer-assisted and
unconditional; Theorem~\ref{thm:intro-D} is proved, and its hypotheses are open. This list is the single place where the limits of the
results are stated.
\begin{enumerate}[label=(\arabic*),leftmargin=*]
\item Nothing here proves RH, \condE, \condS\ or \condU, and Conjecture~\ref{conj:U} itself does not imply RH
(Section~\ref{ssec:intro-pairs}).
\item Only one side of criticality is certified. The inequality $\kstar\ge0$, equivalently $\qmin\le1$, is \condE; we know it only under
RH, and unconditionally we know only $\kstar\ge-2.7112\cdot10^{-3}$, that is, $\qmin<1.017182$. We never claim $\kstar=0$.
\item The hypotheses of Theorem~\ref{thm:intro-D} concern infinitely many $J$, or a limit in $J$, and are open. Only finite-$J$ instances
are certified (Propositions~\ref{thm:finiteJ} and~\ref{prop:exactcone}); convergence of the family is neither assumed nor proved.
\item Uniqueness is proved only near $\zeta$'s support and for integer weights. Theorem~\ref{thm:intro-B}(a) allows any prime measure but no
extra zeros other than $0$; Theorem~\ref{thm:intro-B}(b) allows countably many extra zeros under a summability condition, but no prime mass off
$\BRS$; Theorem~\ref{thm:integer} allows any prime measure, but only integer weights. Extra zeros, when some zero weight is not an integer,
together with prime mass off $\BRS$ or without summability, are open, and so is non-atomic extra zero mass. Corollary~\ref{cor:nearrigid} is quantitative only at
low height.
\item The comparison with $\QQ(\sqrt5)$ and $\QQ(\sqrt{-3})$ concerns two examples, at $\zeta$'s gap and at their natural gaps; it is not a
statement about $L$-functions in general (Remark~\ref{rem:gap}).
\end{enumerate}

\subsection{Related work}

\emph{Explicit formulas and positivity.} The explicit formula goes back to Riemann, von Mangoldt, Guinand \cite{Gui48,Gui49} and Weil
\cite{Wei52,Wei72}. Weil's criterion characterises RH by positivity of a \emph{quadratic} form on test functions of the type $f*\tilde f$
\cite{Wei52,Bom00}; our cone $\Cone$ is defined instead by pointwise positivity of $F$ and $\hF$, and the problem is a linear programme. The
archimedean part of Weil's form is positive on test functions supported in $[2^{-1/2},2^{1/2}]$ whose Fourier transform vanishes at
$\pm\frac i2$; this was proved by Yoshida \cite{Yos92}, and Connes and Consani \cite{CC21a} gave an operator-theoretic strengthening. Connes and
Consani \cite{CC23} exhibited very small eigenvalues of Weil's form on test functions supported in a window $[-L,L]$, and Zhu \cite{Zhu26}
certifies two-sided bounds for them; from the prime side of the explicit formula alone he shows that the normalised smallest eigenvalue
for $L=2$ is at most $3.2\cdot10^{-283}$, a quadratic analogue of Theorem~\ref{thm:intro-C}. Connes \cite{Con26} extremises a restriction
of Weil's form built from the primes $p\le13$ to approximate the first $50$ zeros. None of these works makes a statement about positive
solutions of the explicit formula. Weil's Hermitian form and the explicit formula are also the tools of recent work on a different
question, the proportion of simple zeros on the critical line \cite{AF26,Lam26}; the proof in \cite{AF26} was found by an AI model. An
AI-generated preprint, not peer-reviewed, claims that $\zeta$ and every Dirichlet $L$-function are zero-free in $\Re s>7/8$
\cite{OAI26qrh}; the authors report a Lean formalisation of this statement. It concerns the location of zeros, not positive solutions of
\eqref{eq:intro-EF}, and nothing here uses it. The problem that defines $\kstar$ involves only the archimedean data and the gap, and so does the Gaussian--Laguerre
certificate of Proposition~\ref{thm:kappaOPS}. The certificates of Theorem~\ref{thm:intro-C}(a), in contrast, are built from arithmetic
input: the factor $\Xi^2$, which vanishes at the zeros of $\zeta$, and Hermite interpolation at the prime powers.

\emph{Fourier summation formulas.} Admissible pairs are pairs of positive measures in Fourier duality with fixed archimedean data.
Dyson suggested that a classification of one-dimensional quasicrystals might lead to a proof of RH \cite{Dys09}; such pairs of
sequences go back to Guinand's theory of concordance \cite{Gui59}. Gon\c{c}alves \cite{Gon23} classified Fourier summation formulas,
with $\zeta$'s explicit formula in his class, and his Theorem~5 classifies non-negative crystalline measures with uniformly discrete
support; Gon\c{c}alves and
Vedana \cite{GV25} and Vedana \cite{Ved26} completed the classification on the line and extended it to a strip. These classifications
concern formulas whose summation side is a countable sum of point masses, whereas an admissible pair may have an arbitrary positive prime
measure; so they are related frameworks, which apply directly only to admissible pairs with a discrete prime measure (for instance one
carried by $\BRS$, as in Theorem~\ref{thm:intro-B}(b)). None of these works imposes positivity on both sides with fixed archimedean data.

\emph{Conductor bounds and characterisations of $\zeta$.} The method for discriminants is due to Stark and Odlyzko, with the explicit
formula introduced by Serre and refinements by Poitou \cite{Sta74,Odl76,Ser75,Poi77,Poi77b}; Odlyzko treated Artin conductors \cite{Odl77}
and Mestre conductors of varieties \cite{Mes86}; see \cite{Odl90,DyD80}. In Odlyzko's own setting, the classical cone, a Gaussian--Laguerre
function (Proposition~\ref{thm:kappaOPS}) raises the tabulated $0.997$ to at least $1-6.25\cdot10^{-40}$, and the zero-killing certificate of
Theorem~\ref{thm:intro-C}(a), which also lies in the classical cone, to at least $1-8.98\cdot10^{-1060}$. Both address \cite[Open Problem~2.2]{Odl90} in degree one for the data of $\zeta$. In the literature we
searched we found no optimisation of this bound beyond Odlyzko's one-parameter family, and no claim that its optimum is $1$. A related
extremal problem concerns the lowest zero: among Dedekind zeta functions, $\zeta$ has the highest lowest zero \cite[p.~132]{Odl90};
Miller \cite{Mil02} proved the analogue under GRH for entire $L$-functions of real archimedean type by positivity in the explicit
formula, and Bober et al.\ \cite{Bob15} showed that it fails beyond real archimedean type.
Hamburger's theorem \cite{Ham21,Ham21b} (see \cite{Per17,Bur12}) and its extensions \cite{Boc58,Kno94,Kno00}, the classification of the
degree-one Selberg class \cite{Sel92,CG93,KP99,Sou05}, and the theorems of Beurling \cite{Beu51} and A.~B.~Dixit \cite{DixA20}
characterise $\zeta$ among Dirichlet series, or impose positivity on Dirichlet coefficients. Conjecture~\ref{conj:U} starts instead from
the explicit formula, with arbitrary positive measures on both sides. The closest antecedent we found is Booker's converse theorem for
positive $L$-data \cite[Definition~1.3, Theorem~1.7]{Boo15}; Theorems~\ref{thm:intro-B} and~\ref{thm:integer} are complementary to it, not
generalisations (Sections~\ref{ssec:Gfin} and~\ref{ssec:integer}).

\emph{Interpolation, linear programming and magic functions.} Bondarenko, Radchenko and Seip \cite{BRS} constructed a Fourier
interpolation basis for the zeros of $\zeta$ and the points $\BRS$, after Radchenko and Viazovska \cite{RV19}; Proposition~\ref{thm:support}
is the uniqueness statement implicit in \cite[\S4.4]{BRS}. The linear programming bounds of Cohn and Elkies \cite{CE03} and the magic
functions of \cite{Via17,CKMRV17,CKMRV22} are the model for our use of ``magic function'', as is the uniqueness of the Leech lattice
\cite{CK09}. In that setting near-sharpness is familiar, and our results have direct analogues there: Theorem~\ref{thm:intro-C} corresponds to
the certificate of Cohn and Kumar that the linear programming bound in dimension $24$ is sharp up to a factor $1+1.65\cdot10^{-30}$
\cite{CK09}; the explicit admissible pairs of Propositions~\ref{prop:lowerbound} and~\ref{thm:notcritical} correspond to the dual
certificates with which Cohn and Triantafillou proved non-sharpness in dimensions $12$, $16$, $20$, $28$ and $32$ \cite{CT22}; and Theorem~\ref{thm:intro-D} with
Conjecture~\ref{conj:family} corresponds to the sequences that Cohn and Miller conjectured to converge to magic functions \cite{CM16}. In dimension $2$, an AI-generated preprint, not
peer-reviewed, claims the sharp case of the Cohn--Elkies bound \cite[Conjecture~7.3]{CE03} through a radial Schwartz certificate obtained by
Hermite-cardinal interpolation at a periodic superset of the triangular-lattice shells, with signs certified in interval arithmetic
\cite{OAI26planar}; the authors report a Lean formalisation of the existence of this certificate. The analogue for the explicit formula, an
exact certificate giving \condS, is what Theorem~\ref{thm:intro-D} and Conjecture~\ref{conj:family} ask for; the periodic-node device has
no counterpart for the prime powers. The cone $\Cone$ is a sign-uncertainty cone \cite{BCK10,CG19,CIR24}, and Bourgain, Clozel and Kahane already related such
problems to zeta functions and discriminant bounds \cite[\S4]{BCK10}. Linear programmes over trace and crossing formulas, with near-sharp
bounds attained by special arithmetic objects, appear in the bootstrap \cite{HMR19,KMP24}. Fourier optimisation with the explicit formula
is usually done under GRH \cite{CMS19,CGdL}; our problem is unconditional on the archimedean side and asks about all positive solutions. A
test function whose transform contains $\zeta$ as a factor kills every non-trivial zero, and Guinand's formula then gives an identity over
prime powers alone \cite{RR22}; our certificates use $\Xi^2$ for this purpose, with a cofactor chosen by Hermite interpolation.

\emph{What is new.} As far as our literature search shows, the following are new: the linear programme over all positive solutions of
the explicit formula with $\zeta$'s data, with Conjecture~\ref{conj:U}, its logic and the link between criticality and atomicity
(Theorem~\ref{thm:intro-A}(b)); Theorem~\ref{thm:intro-B} and the Corollary; Theorem~\ref{thm:integer} as an application to the explicit formula (its
analytic ingredients are classical); the zero-killing certificates $\Xi^2H$, which evaluate the dual
objective through $\zeta$'s own explicit formula (Theorem~\ref{thm:intro-C}(a)); the criterion of Theorem~\ref{thm:intro-D}; and the
certified non-criticality of the data of $\QQ(\sqrt5)$ and $\QQ(\sqrt{-3})$ at their natural gaps. Proposition~\ref{thm:support} is essentially contained in
\cite[\S4.4]{BRS} with \cite[Corollary~1.1]{BRS}, and the duality theory of Section~\ref{sec:setting} is standard conic linear programming.

\subsection{Organisation, notation and companion manuscripts}

Section~\ref{sec:setting} sets up admissible pairs and proves duality, the logic and criticality. Section~\ref{sec:uniqueness} proves the
uniqueness results, Section~\ref{sec:family} introduces the zero-killing family and proves Theorem~\ref{thm:intro-D},
Section~\ref{sec:certified} contains the certified bounds, and Section~\ref{sec:conjectures} the conjectures and open problems.
Appendix~\ref{app:certification} describes the certification methods, Appendix~\ref{app:proofs} contains the long proofs, and
Appendix~\ref{app:anc} indexes the ancillary files. A companion manuscript \cite{N1} studies the family of Section~\ref{sec:family} in
detail, and three notes treat extra prime-side atoms \cite{N2}, the critical conductors of other Gamma factors \cite{N3}, and pairs with
integer zero weights \cite{N4}; nothing here depends on them.

We use $\hF(\xi)=\int F(t)e^{-2\pi i\xi t}\dd t$, $\xin n=\log n/(2\pi)$ for real $n>0$, and the multiplicative coordinate $x=e^{2\pi\xi}$, so
that the prime-power node $\xin n$ corresponds to $x=n$ and the gap $[\xin2,\infty)$ to $x\ge2$. $\PP=\{n_1<n_2<\cdots\}$ is the set of prime
powers and $\xiPP=\{\xin n:n\in\PP\}$; $\Zz$ is the set of real $\gamma$ with $\zeta(\frac12+i\gamma)=0$; $\GR(s)=\pi^{-s/2}\Gamma(s/2)$; and
$\psi=\Gamma'/\Gamma$ is the digamma function.
